\documentclass[10pt]{article}
\usepackage[margin=18mm]{geometry}
\usepackage[T1]{fontenc}
\usepackage[utf8]{inputenc}
\usepackage{booktabs}
\usepackage{tabularx}
\usepackage{array}
\usepackage{xcolor}
\usepackage[hidelinks]{hyperref}

\definecolor{navy}{HTML}{17365D}
\definecolor{blue}{HTML}{2F5597}
\definecolor{pale}{HTML}{EAF0F7}
\definecolor{gray}{HTML}{555555}
\hypersetup{colorlinks=true,linkcolor=navy,urlcolor=blue}
\setlength{\parindent}{0pt}
\setlength{\parskip}{4pt}
\renewcommand{\arraystretch}{1.16}

\newcommand{\key}[1]{\textbf{\color{navy}#1}}
\newcommand{\finding}[1]{%
  \begin{center}\colorbox{pale}{\parbox{0.93\linewidth}{\textbf{#1}}}\end{center}}

\begin{document}

\begin{center}
{\LARGE\bfseries\color{navy} Freeway Speed Advice at a Lane Drop}\par
\vspace{2mm}
{\large What the methodology review established, and what remains unproved}\par
\vspace{2mm}
{\small Colleague report \textbar{} 30 September 2026}
\end{center}

\finding{The reviewed freeway task is scientifically legitimate and well anchored in
existing traffic-control and lane-changing literature. The review does not,
however, establish that lane-resolved information is necessary for the control
decision, nor that a lane-aware differentiable model would provide a useful
gradient. The review closes because that boundary is now clear and the evidence
needed to cross it is precisely defined.}

\section{Why this review was needed}

The project asks how much traffic discreteness must be retained to obtain useful
gradients for control. A model-first route would build a lane-aware differentiable
simulator and only later ask whether its extra state changes a meaningful control
decision. This review instead used a task-first anchor: define one freeway control
problem, identify a strong lower-resolution competitor, and require evidence that
the information removed by that competitor changes the decision before proposing
a richer model.

The resulting task is deliberately narrow. Infrastructure supplies the same
upstream speed advice to all lanes approaching a lane reduction. Drivers and
traffic interactions determine actual speeds, lane changes and merging. The
initial controlled setting is homogeneous traffic at a synthetic permanent
three-to-two lane drop. Direct commands to individual vehicles or lane changes
are outside the action space.

The aggregate comparator is intentionally strong. It may know spatial traffic
amount, mean speed or response, queues, cumulative counts and the full declared
predecision history. A lane-aware method must add decision-relevant information;
it cannot win because the baseline was made memoryless or poorly calibrated.

\section{The task that the literature supports}

\begin{tabularx}{\textwidth}{>{\bfseries}p{0.21\textwidth} X X}
\toprule
Task element & Literature foundation & Consequence for this review \\
\midrule
Aggregate VSL baseline & LWR and link-queue studies already regulate lane-drop
traffic with feedback speed limits; first-order VSL sensitivities and gradient
methods also exist \cite{jin2013,monache2016}. & Differentiability and successful
aggregate VSL are prior art. The baseline must be treated as a credible
competitor. \\
\addlinespace
Imperfect response & Macroscopic models and field studies represent incomplete
compliance and the gap between displayed advice and realized speed
\cite{frejo2018,caltrans2018}. & Imperfect response does not by itself require
individual vehicles or lane state. \\
\addlinespace
Lane and merge effects & Lane-resolved CTM, microscopic lane-change models and
field lane-use studies show that mandatory merging, urgency, gaps and lane
distribution can affect traffic near a reduction \cite{guo2020,erdmann2015,knoop2010}. &
Lane, gap and order variables are plausible candidates, but their necessity for
common speed advice remains a hypothesis. \\
\addlinespace
Richer reference evaluation & Existing work compares aggregate controllers with
microscopic simulation, including SUMO-based evaluation \cite{guo2020,chen2025}. &
A richer evaluator can test transfer, but richer trajectories alone do not prove
that added information improves the control decision. \\
\addlinespace
Temporary-event recovery & Incident guidance treats the period after lane
reopening, while queues dissipate, as a distinct operational phase
\cite{fhwa2008}. & A later temporary-closure review is credible, but it is a new
scenario rather than missing evidence for the permanent-drop review. \\
\bottomrule
\end{tabularx}

No reviewed paper supplies the complete target comparison: the same lane-uniform
advice, the same strong aggregate information, a reachable hidden lane/gap/order
contrast, complete waiting and travel accounting, a minimal differentiable
refinement, and gradient-guided improvement in an independent richer evaluator.

\section{The three central questions and their answers}

\subsection{1. What information could an aggregate simulator hide?}

The physically relevant candidates are near-drop lane allocation, continuing-lane
balance, local gaps, queue or packet order, leaders and lane-change memory. Lane
occupancy is best understood as a marker for this bundle, rather than an already
proved sufficient state.

The review obtained one conditional analytical result. In a homogeneous
free-flow preparation, conservative transfers between lanes cancel when traffic
mass is summed. Different lane allocations can therefore share the same aggregate
mass history under the stated assumptions. This demonstrates how information can
be hidden without rounding vehicle mass or disabling conservation.

It does not demonstrate a native near-drop case. Once lane-dependent speeds,
receiving constraints and merging affect aggregate measurements, the full history
may reveal the preparation. The review designed a bounded future test for this
existence question, but did not run it.

\subsection{2. Can hidden information change useful common speed advice?}

The literature establishes plausible mechanisms: speed affects arrivals, gap
formation, lane distribution and merge interactions. It does not isolate a case
where two situations with the same strong aggregate history require different
lane-uniform advice after all waiting traffic is counted.

The analytical review also found an important obstruction. If advice changes
only when traffic arrives, while bottleneck service is fixed, advice-independent
and work-conserving, sending traffic earlier cannot be worse under complete delay
accounting. The highest feasible advice remains a common best action. Assigning
different fixed capacities to hidden lane states therefore does not establish a
reason to slow traffic.

A successful case must contain an endogenous coupling: common advice must change
merging, receiving or service enough to offset its added travel delay, and this
effect must depend on information unavailable in the aggregate history. No
reviewed source establishes that full chain for this task.

\subsection{3. Can the missing information provide a useful gradient?}

The required evidence path is clear:

\begin{enumerate}
\item establish a reachable collision in the declared aggregate information;
\item establish different useful common-advice decisions;
\item identify the smallest retained state that explains the difference;
\item compare automatic or analytic derivatives with finite differences away
      from event boundaries; and
\item replay small gradient-guided perturbations in an independent richer
      evaluator.
\end{enumerate}

The review completed the methodological definition of this path, not the traffic
instance or derivative evidence. A differentiable simulator, a finite derivative
or a lower internal objective would not by itself establish gradient usefulness.

\section{What is established, and where the chain stops}

\begin{tabularx}{\textwidth}{>{\bfseries}p{0.27\textwidth} X}
\toprule
Established result & Scientific meaning \\
\midrule
Strong aggregate comparator & Aggregate response/history models are a serious
alternative and must not be weakened to create a gap. \\
Conditional hiding construction & Conservative lane redistribution can be
invisible in aggregate free-flow history under explicit assumptions. \\
Fixed-service obstruction & Hidden lane state plus different fixed capacity is
insufficient when advice cannot change service and faster arrival dominates. \\
Mechanism foundation & Gap, urgency, lane redistribution and native merging are
plausible paths, without a guaranteed beneficial sign. \\
Reachable-preparation repair & Later decisions and a localized balanced contrast
avoid the most obvious empty-near-zone and unequal-insertion confounds. \\
Future exact-prefix protocol & A frozen dual-resolution design can test exact
sampled-history collision for one bounded synthetic preparation bank. \\
\bottomrule
\end{tabularx}

\vspace{2mm}
\key{The chain stops before the decisive link:} no source or proof combines a
reachable strong-history match, consequential hidden state and different useful
common advice. Therefore no minimum state and no task-specific useful gradient
can yet be claimed.

\section{Why the remaining gaps cannot be closed by more static review}

\begin{itemize}
\item \key{The information collision is existential.} Literature can justify
      mechanisms and observer choices, but it cannot certify that the frozen
      native preparations produce exactly equal histories.
\item \key{The control effect is confounded in close predecessors.} Existing
      lane-aware studies often change the actuator, objective, optimizer or
      observability together with model resolution.
\item \key{The minimum state depends on a positive case.} Lane mass, gaps, order
      and lane-change memory cannot be ranked by necessity before a
      decision-relevant contrast exists.
\item \key{The derivative target depends on the adopted mechanism.} Exact
      propagation through a conservative smooth branch differs from sensitivity
      at a lane-change or reopening event.
\item \key{A complete gradient claim requires construction and evaluation.}
      Internal derivative checks and independent downstream usefulness require a
      candidate model and controlled perturbations, neither of which belongs to
      a methodology-only literature review.
\end{itemize}

These are evidence boundaries, not clerical omissions. Inventing an advice-to-
service equation or installing a microscopic simulator would start a new study.

\section{Conclusion}

The current review closes with conditional findings. The formulated permanent
lane-drop VSL task is defensible and literature-grounded, and the evidence needed
for a selective mesoscopic contribution is explicit. The central hypothesis
remains unresolved rather than refuted.

The safe claim is therefore methodological: the literature supports the task,
the strong aggregate baseline and the physical relevance of lane/gap/order
effects, while the review identifies the exact evidence needed before selective
mesoscopic state or useful gradients can be claimed. Establishing that evidence
requires a new analytical or empirical study rather than further reinterpretation
of the current literature.

\small
\begin{thebibliography}{99}
\bibitem{jin2013} W.-L. Jin and H. Jin, ``Analysis and design of variable speed
limit control for a lane-drop bottleneck,'' 2013. Available:
\url{https://arxiv.org/abs/1310.2658}.
\bibitem{monache2016} M. L. Delle Monache, B. Piccoli, and F. Rossi,
``Traffic regulation via controlled speed limit,'' 2016. Available:
\url{https://arxiv.org/abs/1603.04785}.
\bibitem{frejo2018} J. R. D. Frejo et al., ``Macroscopic modeling and control of
variable speed limits with non-compliance,'' \emph{IFAC-PapersOnLine}, 2018,
doi: 10.1016/j.ifacol.2018.07.056.
\bibitem{caltrans2018} California PATH, ``Field test of variable speed advisory
system,'' Caltrans Final Report CA18-2447, 2018. Available:
\url{https://dot.ca.gov/-/media/dot-media/programs/research-innovation-system-information/documents/final-reports/ca18-2447-finalreport-a11y.pdf}.
\bibitem{guo2020} Y. Guo et al., ``Integrated variable speed limits and
lane-changing control for freeway lane-drop bottlenecks,'' \emph{IEEE Access},
vol. 8, pp. 54710--54721, 2020, doi: 10.1109/ACCESS.2020.2981658.
\bibitem{erdmann2015} J. Erdmann, ``SUMO's lane-changing model,'' in
\emph{Modeling Mobility with Open Data}, 2015. Available:
\url{https://elib.dlr.de/102254/1/Springer-SUMOs_Lane_changing_model.pdf}.
\bibitem{knoop2010} V. L. Knoop et al., ``Lane distribution of traffic near
merging zones: influence of variable speed limits,'' IEEE ITSC, 2010,
doi: 10.1109/ITSC.2010.5625034.
\bibitem{chen2025} Y. Chen, G. Cheng, F. Meng, and W. Wang, ``Variable speed
limit control method of freeway merging area in a car-truck mixed heterogeneous
traffic environment,'' \emph{Physica A}, vol. 676, 130876, 2025,
doi: 10.1016/j.physa.2025.130876.
\bibitem{fhwa2008} Federal Highway Administration, ``Information sharing for
traffic incident management,'' FHWA-HOP-08-059, 2008. Available:
\url{https://ops.fhwa.dot.gov/publications/fhwahop08059/30.htm}.
\end{thebibliography}

\end{document}
